\chapter{Cache Coherence and Performance}
\label{chap:cacheperformance}

Modern CPUs are fast; \highlight{DRAM} is slow. Caches hide that gap.
\par\noindent\textbf{Interview tip:} explain why two threads touching 
\highlight{nearby} memory can destroy performance (\highlight{false sharing}), 
how coherence protocols move ownership of cache lines, and how to 
\highlight{measure} misses with \texttt{perf}.

\section{CPU Cache Hierarchy}
\label{sec:cpu-cache-hierarchy}
Caches store recently used data and instructions in small, fast SRAM. 
Capacity grows and latency worsens as you go down the hierarchy. Typical 
orders of magnitude on a server CPU (rough; learn the story, not exact ns):
\begin{itemize}
    \item L1: $\sim$1--4 cycles, tens of KiB per core.
    \item L2: $\sim$10+ cycles, hundreds of KiB per core.
    \item L3: tens of cycles, shared across cores on a socket, MBs.
    \item DRAM: hundreds of cycles.
\end{itemize}
A \highlight{cache line} is the transfer unit (commonly 64 bytes on current x86-64 CPUs). 
The hardware loads/invalidates whole lines, not single \texttt{int}s.

\subsection{L1}
\label{subsec:cache-l1}
\highlight{L1} is split into instruction (L1i) and data (L1d) on most cores. 
It is private to the core, tiny, and the first place a hot trading loop 
should keep its working set. A miss in L1 often hits L2 next.
\newline
\textbf{Noise removed} when the hot working set fits in L1: DRAM-scale stalls on every load. 
Cost of missing: a few to tens of cycles per access --- still ``local,'' but enough to matter 
when the loop is already nanosecond-sensitive. Instruction-cache misses (cold code paths, 
giant binaries) add the same class of jitter on branches you rarely take until a rare event.

\subsection{L2}
\label{subsec:cache-l2}
\highlight{L2} is larger and still typically \highlight{per-core} (or lightly shared). 
It absorbs what does not fit in L1. Pinning a thread to one core 
(Chapter~\ref{chap:linuxprocesses}) helps keep L1/L2 \highlight{warm}.
\newline
\textbf{Noise removed} by keeping the mid-size working set in L2 after L1 overflows: 
repeated trips to the shared LLC or DRAM. Cost of missing L2: typically tens of cycles 
plus contention if the line must come from another core or L3. Migration undoes L1/L2 warmth 
--- that is why pinning is a cache story, not only a scheduler story.

\subsection{L3}
\label{subsec:cache-l3}
\highlight{L3} (LLC) is usually \highlight{shared} by cores on the same socket. 
It helps when threads on the same socket share read-mostly data. 
It does \highlight{not} remove false sharing: two cores writing the same line 
still bounce ownership through the coherence protocol.
\newline
\textbf{Noise removed} when shared read-mostly state stays in L3 instead of DRAM: 
hundreds of cycles per miss and interconnect pressure. Cost: other cores on the socket 
compete for the same LLC capacity --- noisy neighbors thrash your lines even when your 
thread is pinned. Remote DRAM (Chapter~\ref{chap:numa}) is the next cliff after L3.
\par\noindent\textbf{Interview tip:} frame each level by \emph{working set that fits} and 
\emph{what you pay when it does not} --- L1 (cycles), L2 (tens), L3 (tens--hundreds), DRAM (hundreds+).

\section{Cache Coherence}
\label{sec:cache-coherence}
When multiple cores cache the same physical line, they must agree on its 
contents. \highlight{Cache coherence} protocols maintain that agreement. 
Without them, core A could keep a stale copy while core B wrote a new value.

\subsection{MESI}
\label{subsec:mesi}
See~\cite{intel-mesif}.
\newline
\highlight{MESI} is a classic write-invalidate protocol. Each cached line is in one of:
\begin{itemize}
    \item \highlightbf{M (Modified)} --- this core has the only copy; it is dirty (differs from DRAM).
    \item \highlightbf{E (Exclusive)} --- this core has the only copy; clean.
    \item \highlightbf{S (Shared)} --- several cores may hold a read-only copy.
    \item \highlightbf{I (Invalid)} --- not present / must not be used.
\end{itemize}
A write typically requires \highlight{Exclusive} or \highlight{Modified} ownership: 
other cores' copies are invalidated first.

\begin{figure}[H]
\centering
% MESI states (externalized). Legend is a 2-column tabular so 2/4/6 line up.
\begin{minipage}{\textwidth}
\centering
\tikzfig{mesi_protocol}
\begin{tikzpicture}[
    font=\sffamily\footnotesize,
    state/.style={
        draw=black,
        thick,
        circle,
        minimum size=1.05cm,
        align=center,
        inner sep=1pt
    },
    arr/.style={-{Stealth[length=2mm]}, thick},
    lab/.style={font=\sffamily\scriptsize, text=oblivioncomment, align=center},
    tag/.style={font=\sffamily\tiny\bfseries, fill=white, inner sep=1.5pt}
]
    \node[state, fill=myblueiii]  (M) at (0,2.6) {\textbf{M}};
    \node[state, fill=whitesmoke] (E) at (4.0,2.6) {\textbf{E}};
    \node[state, fill=whitesmoke] (S) at (4.0,0) {\textbf{S}};
    \node[state, fill=lightgray]  (I) at (0,0) {\textbf{I}};

    \node[lab, above=0.28cm of M] {Modified\\{\scriptsize dirty, unique}};
    \node[lab, above=0.28cm of E] {Exclusive\\{\scriptsize clean, unique}};
    \node[lab, below=0.36cm of S] {Shared\\{\scriptsize clean, copies OK}};
    \node[lab, below=0.36cm of I] {Invalid};

    \draw[arr] (I) -- node[pos=0.5, above left=1pt and 2pt, tag] {1} (E);
    \draw[arr] (E) -- node[pos=0.5, above=5pt, tag] {2} (M);
    \draw[arr] (E) -- node[pos=0.5, right=5pt, tag] {3} (S);
    \draw[arr] (S) -- node[pos=0.5, below=5pt, tag] {4} (I);
    \draw[arr] (M) -- node[pos=0.5, left=5pt, tag] {5} (I);
    % 6: toward S (down/right of the crossing), still on the dashed bend
    \draw[arr, dashed] (S) to[bend left=22]
        node[pos=0.28, below right=1pt and 2pt, tag] {6} (M);

\end{tikzpicture}

\vspace{0.35em}
{\sffamily\scriptsize\color{oblivioncomment}%
\begin{tabular}{@{}r@{\,:\,}p{5.4cm}@{\hspace{1.2em}}r@{\,:\,}p{5.4cm}@{}}
1 & read miss / allocate (I$\rightarrow$E) &
2 & local write (E$\rightarrow$M) \\
3 & other core reads (E$\rightarrow$S) &
4 & remote write / invalidate (S$\rightarrow$I) \\
5 & writeback + invalidate (M$\rightarrow$I) &
6 & write while Shared (dashed; invalidate others) \\
\end{tabular}}
\end{minipage}

\caption{MESI states (sketch). Arrows are common transitions, not the full protocol.}
\label{fig:mesi-protocol}
\end{figure}

\toclessnparagraph{Walkthrough (two cores, one line)}
\begin{enumerate}
    \item Both start \textbf{I}. Core~A reads $\Rightarrow$ line arrives in \textbf{E} on A.
    \item Core~B reads the same address $\Rightarrow$ both hold \textbf{S} (A drops from E to S).
    \item Core~A writes $\Rightarrow$ it must \highlight{invalidate} B's copy (B $\rightarrow$ \textbf{I}); 
          A's line becomes \textbf{M}.
    \item Core~B reads again $\Rightarrow$ miss; it may force A to write back, then both can end in 
          \textbf{S} (or B gets a fresh copy and A is demoted). Exact bus messages vary by CPU; 
          the interview point is \highlight{invalidation + refetch cost}.
\end{enumerate}
That bounce is what you feel as \highlight{coherence traffic}. False sharing 
(Section~\ref{sec:false-sharing}) is the same mechanism when the ``shared'' line 
holds unrelated variables.

\subsubsection{Cache Invalidation}
\label{subsec:cache-invalidation}
When core A writes a line held as Shared elsewhere, it sends 
\highlight{invalidations}. Other cores mark their copies Invalid. 
Their next access \highlight{misses} and must refetch. 
Frequent invalidations across cores = \highlight{coherence traffic} and latency spikes.

\par\noindent\textbf{Interview tip:} name the four letters, then walk one write that 
turns Shared copies into Invalid --- that is the link to false sharing.

\section{False Sharing}
\label{sec:false-sharing}
\par\noindent\textbf{Interview tip:} false sharing is one of the highest-yield topics in this chapter.

\subsection{What It Is}
\label{subsec:false-sharing-definition}
\highlight{False sharing} happens when two threads update \highlight{different variables} 
that live on the \highlight{same cache line}. Logically there is no shared data race 
on one variable; physically the line is shared and bounced between cores.

\begin{lstlisting}[language=C++]
#include <atomic>
#include <thread>

struct SharedCounters {
    std::atomic<long> a{0};
    std::atomic<long> b{0};  // may share a cache line with a
};

void false_sharing_demo(SharedCounters& c)
{
    std::thread t1([&] {
        for (int i = 0; i < 100000000; ++i) c.a.fetch_add(1, std::memory_order_relaxed);
    });
    std::thread t2([&] {
        for (int i = 0; i < 100000000; ++i) c.b.fetch_add(1, std::memory_order_relaxed);
    });
    t1.join();
    t2.join();
}
\end{lstlisting}

\subsection{Why It Slows Things Down}
\label{subsec:false-sharing-why-slow}
Each store may force the other core to \highlight{invalidate} and reload the line. 
You pay coherence traffic even though \texttt{a} and \texttt{b} are independent. 
Symptoms: poor scaling when adding threads, high cache-miss / remote-hit counters, 
time spent in what looks like ``simple atomics.''

\subsection{Cache-Line Ping-Pong}
\label{subsec:cache-line-ping-pong}
The line oscillates (``\highlight{ping-pong}'') between cores: M on core 0, then M on core 1, 
and so on. Hot SPSC indices without padding show the same pathology 
(Section~\ref{subsec:spsc-cache-line-layout}).

\subsection{Solutions}
\label{subsec:false-sharing-solutions}
\begin{itemize}
    \item \highlightbf{Pad / align} so hot fields sit on different lines. 
          Examples below use \texttt{alignas(64)} because 64 bytes is the usual 
          line size on x86-64 and keeps listings readable. Prefer the portable 
          C++17 constant \texttt{std::hardware\char`_destructive\char`_interference\char`_size} 
          (\texttt{<new>}) in production code --- it is the destructive interference size 
          (padding to \highlight{avoid} sharing a line). The related 
          \texttt{std::hardware\char`_constructive\char`_interference\char`_size} is for 
          packing data that \highlight{should} share a line.
\end{itemize}
\begin{lstlisting}[language=C++]
#include <atomic>

struct PaddedCounters {
    alignas(64) std::atomic<long> a{0};
    alignas(64) std::atomic<long> b{0};
};
\end{lstlisting}
\begin{itemize}
    \item Give each thread its \highlight{own} counter (or shard), then reduce infrequently.
    \item Keep writers on the \highlight{same core} if they must share a line (often worse than padding).
    \item Verify with \texttt{perf} (Section~\ref{sec:cache-profiling}).
\end{itemize}

\section{Profiling}
\label{sec:cache-profiling}
Guessing is expensive; \highlight{measure}. On Linux, \texttt{perf} is the standard first tool.
\newline
Decision narrative: start wide (\texttt{perf stat}), decide whether the bottleneck looks 
cache-like, then zoom to code (\texttt{record}/\texttt{annotate}), then classify the miss 
pattern (sharing vs capacity vs NUMA). Do not jump to padding or huge pages until the 
counters and samples point there.

\subsection{\texttt{perf stat}}
\label{subsec:perf-stat}
High-level counters for a whole run:
\begin{lstlisting}[language=bash]
perf stat -e cycles,instructions,cache-misses,cache-references ./my_bench
\end{lstlisting}
Look at \texttt{cache-misses} relative to \texttt{cache-references}, and 
instructions per cycle (IPC). A sudden miss rate jump after a ``harmless'' 
structural change often points at layout / false sharing.
\newline
\textbf{Noise removed} by this step: debating micro-optimizations while the miss rate 
was never the problem --- or while IPC was already healthy and the stall was I/O or locks.

\subsection{\texttt{perf record} and \texttt{perf annotate}}
\label{subsec:perf-record-annotate}
\begin{lstlisting}[language=bash]
perf record -e cache-misses -g ./my_bench
perf report
perf annotate
\end{lstlisting}
\texttt{record} samples where misses occur; \texttt{annotate} maps them to 
instructions / source. Use this when \texttt{stat} says misses are high but 
you do not know \highlight{which} line of code.
\newline
If samples cluster on two atomics a few bytes apart, suspect false sharing 
(Section~\ref{sec:false-sharing}). If they fan out across a large array scan, suspect 
capacity / layout. If they spike only under multi-socket load, open the NUMA chapter.

\subsection{Cache Miss Events}
\label{subsec:cache-miss-events}
Useful event names vary by CPU; common ones include 
\texttt{cache-misses}, \texttt{cache-references}, and more specific L1/LLC 
miss events from \texttt{perf list}. On Intel, also watch for 
remote DRAM / cross-socket accesses when NUMA is in play 
(Chapter~\ref{chap:numa}).
\newline
Prefer the most specific event that answers your question: L1d misses for a tight loop, 
LLC misses for working-set overflow, remote-hit / DRAM events when pinning or allocation 
locality is in doubt. Generic \texttt{cache-misses} alone does not tell you \emph{which} 
level failed.

\subsection{Interpreting Results}
\label{subsec:profiling-interpretation}
\begin{itemize}
    \item High misses + poor scaling with thread count $\rightarrow$ suspect sharing / false sharing.
    \item High misses in a single-threaded scan $\rightarrow$ suspect layout, random access, or working set $>$ cache.
    \item Good IPC and low misses $\rightarrow$ look elsewhere (syscalls, locks, I/O).
\end{itemize}
\par\noindent\textbf{Interview tip:} say the sequence --- \texttt{stat} to see it, \texttt{record} to find it, 
event type / scaling behavior to name the cause --- not just ``we ran \texttt{perf}.''

\subsection{Detecting False Sharing}
\label{subsec:profiling-false-sharing}
Practical recipe:
\begin{enumerate}
    \item Benchmark the contended structure; note throughput / \texttt{perf stat}.
    \item Apply \texttt{alignas(64)} (or pad) between hot fields; see the note in 
          Section~\ref{subsec:false-sharing-solutions} on 
          \texttt{std::hardware\char`_destructive\char`_interference\char`_size}.
    \item Re-run: large improvement with the same logic $\Rightarrow$ false sharing was likely.
\end{enumerate}
Some toolchains expose \texttt{perf c2c} (cache-to-cache) analysis on supported platforms --- 
useful on real servers when available.

\section{Cache Locality}
\label{sec:cache-locality}
Locality is why contiguous, predictable access patterns beat pointer-chasing.

\subsection{Spatial Locality}
\label{subsec:spatial-locality}
\highlight{Spatial locality}: using address $X$ makes nearby addresses ($X+64$, \ldots) cheap 
because they may already be in the same line or prefetched. 
Prefer \highlight{arrays of structs} laid out for the hot scan order; avoid scattering 
hot fields across the heap.

\subsection{Temporal Locality}
\label{subsec:temporal-locality}
\highlight{Temporal locality}: reusing the same data soon keeps it in L1/L2. 
Tight loops over a small working set win; thrashing a huge working set 
forces DRAM traffic. Pinning (Chapter~\ref{chap:linuxprocesses}) protects 
temporal locality from migration.

\section{TLB}
\label{sec:tlb}
See Section~\ref{sec:address-space} for virtual address space and pages.
\newline
Caches hide DRAM latency for \emph{data}; the \highlight{TLB} hides page-table walk latency 
for \emph{addresses}. A hot path can look cache-friendly in L1/L2 counters and still stall 
on translations if it touches too many pages.

\subsection{Translation Lookaside Buffer}
\label{subsec:translation-lookaside-buffer}
The \highlight{TLB} caches virtual$\rightarrow$physical translations. 
Every load/store needs a translation; a TLB hit is cheap compared to a page-table walk.
\newline
\textbf{Noise removed} by a warm TLB: multi-level page-table walks (and the cache misses those 
walks themselves cause) on every new page the hot loop touches.

\subsection{TLB Misses}
\label{subsec:tlb-misses}
A \highlight{TLB miss} walks page tables (and may incur further cache misses). 
Large, sparsely accessed working sets or many small pages increase TLB pressure. 
With 4~KiB pages, a few dozen MiB of randomly touched memory already needs thousands of 
TLB entries --- more than a typical L1 TLB holds --- so misses become steady noise, not a 
one-time startup cost.
\newline
\textbf{Noise removed} by \highlight{huge pages} (2~MiB / 1~GiB where supported): far fewer 
translations for the same physical span, so the TLB covers a much larger working set. 
Cost: less flexible allocation, setup/\texttt{mmap} flags or transparent huge pages to get right, 
and wasted RAM if you oversize. Still pair with compact layout --- huge pages do not fix 
pointer-chasing across a sparse heap. Page faults and demand paging are covered in 
Section~\ref{subsec:page-faults}.
\par\noindent\textbf{Interview tip:} TLB miss $\neq$ cache miss; huge pages buy fewer translations, 
not faster DRAM --- measure both when the working set is large.
